\session{1}{10/3}{1限}{1}{ガイダンス／AI時代の経営情報}

\goals{
  \item 講義の進め方・評価の考え方・AI活用の方針（AI Policy）を理解する。
  \item 「経営情報」とは何か、情報が経営資源と呼ばれる理由を説明できる。
  \item 生成AIの普及で経営において「変わること」と「変わらないこと」を自分の言葉で整理する。
}

\guidesec{解説}

\topic{講義の進め方}
本講義はオンラインで全15回、4つのモジュールで構成されます（Module 01 経営と情報の基礎、02 企業と情報システム、
03 データと意思決定、04 デジタル時代の経営）。毎回、前半で概念の解説、後半で生成AIを実際に使う演習を行います。
出席確認は毎回のAIクイズ（5問）で行い、全問正解で表示される出席カードを Teams の Class Notebook に提出します。
講義の最後（第14回）には、各自が選んだ企業・組織へのAI活用提案を発表し、相互評価を行います。
評価方法の詳細は講義用 Teams の履修ガイドで確認してください。評価では、AI活用の透明性（どこでどうAIを使ったかを
明示しているか）も見ます。

\topic{経営情報とは何か}
企業が事業を行うために必要な資源は、伝統的にヒト・モノ・カネの3つで説明されてきました。これに「情報」を加えて
4つの経営資源とする見方が、1980年代以降に定着しました。情報が資源であるということは、情報もまた、
集め、整え、活用し、守るという「経営」の対象になるということです。

「経営情報」という分野（英語では Management Information Systems、MIS）は、次の3つの要素の関係を扱います。
\begin{itemize}
  \item \textbf{情報}　経営の意思決定に使われるデータ・情報・知識。
  \item \textbf{情報システム}　情報を集め、処理し、伝える技術のしくみ（ソフトウェア、ハードウェア、ネットワーク、そしてAI）。
  \item \textbf{組織と人}　情報を使って判断し、行動する経営者・管理者・従業員と、その組織。
\end{itemize}
技術だけを学ぶ科目ではなく、技術が経営の意思決定と組織にどう結びつくかを学ぶ科目だ、という点を押さえてください。

\topic{生成AIの普及で変わること、変わらないこと}
2022年末以降、文章・画像・コードなどを生成する生成AIが急速に普及し、情報を「集める・読み解く・書く・提案する」
という作業の多くをAIが担えるようになりました。これにより変わることと変わらないことを整理します。
\par\noindent\begin{gtabh}{colspec={X[l] X[l]}}
  変わること & 変わらないこと \\
  情報収集・要約・下書きの速さと量（人の数十倍） & 何を目的に、何を判断するかを決めるのは人 \\
  情報の「作成」コストの低下（文章・画像・コード） & 判断の結果に責任を負うのは人・組織 \\
  専門知識へのアクセスの平準化（誰でも一定水準の答えが得られる） & 出力が正しいかを見抜くには専門知識が必要 \\
  定型業務の自動化の範囲（判断を伴う業務にも広がる） & 情報の質（正確さ・適時性・関連性）が意思決定の質を決める \\
  求められるスキル（問いを立てる力、評価する力） & 組織の目的・価値観・倫理 \\
\end{gtabh}

\topic{本講義の AI Policy}
AIとの向き合い方として「使う・検証する・区別する」の3つを掲げます。AIを使わないことは安全策ではありません。
使い方を学ばなければ、組織内で無許可のAI利用（シャドーAI、第13回）が広がり、競争力も落ちます。
一方で、AIの出力を鵜呑みにすると、もっともらしい誤り（ハルシネーション、第3回）を経営判断に持ち込むことになります。
だから「使い、検証し、AIの寄与と自分の判断を区別して示す」のです。使用するAIは大学が提供する Copilot です。

\topic{キーワード}
\par\noindent\begin{keywords}{}
  経営資源 & 事業に必要な資源。ヒト・モノ・カネ・情報の4つで整理される。 \\
  経営情報（MIS） & 情報・情報システム・組織と人の3要素の関係を扱う分野。 \\
  生成AI & 学習したパターンをもとに文章・画像・音声などの新しいコンテンツを生成するAI。 \\
  AI Policy & 本講義のAIとの向き合い方。使う・検証する・区別する。 \\
  AI活用記録 & 提出物に添える、AIの使用場面・プロンプト・出力の扱い・検証内容の記録。 \\
\end{keywords}

\guidesec{演習課題}

\exercise{1}{AI演習：「経営情報とは何か」をAIに尋ね、どこまで信じてよいかを検討する}
\instr{Copilot に次のプロンプトを入力し、回答を保存してください。そのうえで、回答の中の (1) 固有名詞、(2) 数値、(3) 断定的な表現を抜き出し、それぞれ「確かめられる／確かめられない／誤っている」に分類してください。}
\begin{promptbox}[経営情報の説明を求める]
あなたは経営学を専門とする大学教員です。大学2年生向けに「経営情報とは何か」を、歴史的な背景、主な研究テーマ、代表的な研究者や文献を含めて400字程度で説明してください。
\end{promptbox}
\instr{次に「今の回答のうち、根拠が弱い部分や、確認が必要な部分はどこですか」と続けて尋ね、AI自身の答えと自分の分類を比べてください。}

\exercise{2}{変わること・変わらないことを自分の経験で言い換える}
\instr{アルバイト、サークル、家庭、就職活動など、自分に身近な組織や活動を1つ選び、生成AIの普及によって「変わること」を3つ、「変わらないこと」を3つ、それぞれ理由とともに書いてください。}

\exercise{3}{AI活用記録を書いてみる}
\instr{演習1でのAIの利用について、付録の形式でAI活用記録を書いてください。「出力をどう扱ったか」と「検証したこと」は具体的に書きます。}

\guidesec{模範解答}

\begin{answer}{演習1}
回答例（Copilot の出力の一部）：「経営情報システム（MIS）は1960年代にアメリカで生まれ、ゴードン・デイビスが1974年に体系化した。
日本では1990年代に経営情報学会が設立され、現在は年間約500本の論文が発表されている。」

分類の例：
\begin{itemize}
  \item 固有名詞「ゴードン・デイビス」「1974年」：デイビス（Gordon B. Davis）の \textit{Management Information Systems}
        は1974年刊行で、確かめられる（大学図書館の蔵書検索で確認）。
  \item 固有名詞「経営情報学会」：実在するが、設立は1992年であり「1990年代」は正しい。細部は学会の公式サイトで確認。
  \item 数値「年間約500本」：出典がなく、確かめられない。使わない。
  \item 断定「1960年代にアメリカで生まれ」：おおむね妥当だが、起点の取り方で説が分かれるため「〜とされる」と書く。
\end{itemize}
AIに根拠の弱い部分を尋ねると、数値について「正確な統計ではない」と答えることが多い一方、人名や年号の誤りは
自分では指摘できないことがあります。「AIに確認させる」だけでは検証にならず、独立した情報源で確かめる必要がある、
というのがこの演習の結論です。
\end{answer}

\begin{answer}{演習2}
例：飲食店のアルバイト
\begin{itemize}
  \item 変わること：(1) シフト表や発注の下書きをAIが作れるので、店長の事務時間が減る。(2) 外国人客への対応で翻訳AIを使い、
        語学力の差が小さくなる。(3) 新人向けマニュアルを写真や動画からAIが生成でき、教育の準備が速くなる。
  \item 変わらないこと：(1) 混雑時に誰をどこに配置するかの判断は店長が現場を見て決める。(2) 食材の品質や衛生の責任は店が負う。
        (3) 常連客との関係づくりは人の接客で決まる。
\end{itemize}
評価の観点：「変わること」が作業の速さ・量・コストに関わり、「変わらないこと」が目的・判断・責任・関係に関わる、
という対応が見えていれば十分です。
\end{answer}

\begin{answer}{演習3}
\par\noindent\begin{gtab}{colspec={Q[l,wd=34mm] X[l]}}
  使用したツール & Copilot（大学提供） \\
  使用した場面 & 「経営情報とは何か」の説明文の生成と、根拠の弱い部分の自己点検 \\
  主なプロンプト & 「あなたは経営学を専門とする大学教員です。…400字程度で説明してください」「根拠が弱い部分はどこですか」 \\
  出力をどう扱ったか & 分野の全体像の把握には採用。論文数の数値は出典がないため不採用。人名・年号は確認のうえ採用。 \\
  検証したこと & デイビスの著書の刊行年を大学図書館の蔵書検索で確認。経営情報学会の設立年を学会公式サイトで確認。 \\
\end{gtab}
\end{answer}

\guidesec{出席確認クイズの解説}
講義サイトの出席確認ページ（第1回）の5問です。正解の選択肢に［正解］を付けています。
% 出席確認クイズ 第01回　ガイダンス／AI時代の経営情報
%   attendance/attendance01.html から生成（解説は本ガイド用に加筆）。

\quizq{1}{「経営情報」という科目で扱う「情報」の位置づけとして最も適切なものはどれでしょうか?}%
  {社内メールの総称}%
  {\correct{ヒト・モノ・カネと並ぶ経営資源のひとつ}}%
  {会計帳簿の別名}%
  {通信回線の速度のこと}%
  {正解は (b)。経営資源はヒト・モノ・カネに「情報」を加えた4つで整理されることが多く、本講義はこの第4の資源としての情報を扱います。(a)(c)(d) はいずれも情報の一側面や伝達手段にすぎず、経営資源としての位置づけを表していません。情報が資源であるなら、集める・整える・活用する・守るという「経営」の対象になる、という発想が出発点です。}

\quizq{2}{生成AI(Generative AI)の説明として適切なものはどれでしょうか?}%
  {家電を遠隔操作するための技術}%
  {与えられたデータを分類するだけのAI}%
  {コンピュータウイルスを検知するソフト}%
  {\correct{文章・画像・音声などの新しいコンテンツを生成するAI}}%
  {正解は (d)。生成AIは、学習したパターンをもとに文章・画像・音声などを新たに「生成」する点が特徴です。(b) の分類や判別は従来型（識別系）AIの典型で、生成AIとは区別されます。(a)(c) はAIの応用先や別分野の技術であって、定義ではありません。第3回でしくみを詳しく扱います。}

\quizq{3}{AIの出力を経営判断に使うときの基本姿勢として適切なものはどれでしょうか?}%
  {AIの出力は常に正しいので、そのまま採用する}%
  {\correct{出力の根拠と限界を確認したうえで、人が判断する}}%
  {回答が長いほど信頼できると考える}%
  {AIは使わないのがもっとも安全である}%
  {正解は (b)。AI Policy の「検証する」にあたる姿勢です。(a) はハルシネーション（第3回）の危険を無視しています。(c) は回答の長さと正確さに関係がないため誤りです。(d) は「使う」という本講義の方針に反し、AIを使わないこと自体がリスク（競争力の低下、シャドーAIの発生）になる場合もあります。}

\quizq{4}{「AI」の正式名称はどれでしょうか?}%
  {\correct{Artificial Intelligence}}%
  {Automatic Information}%
  {Advanced Internet}%
  {Analog Interface}%
  {正解は (a)。AI は Artificial Intelligence（人工知能）の略です。他の選択肢は実在しない語の組み合わせです。語源を押さえると、AIは「人工の」知能であり、万能でも魔法でもなく、設計と学習データに依存した道具だと意識しやすくなります。}

\quizq{5}{本講義の「AI Policy(AIとの向き合い方)」として掲げられて\underline{いない}ものはどれでしょうか?}%
  {AIの寄与と自分の判断を区別して示す}%
  {毎回の演習でAIを実際に使う}%
  {\correct{AIの利用を隠す}}%
  {出力を検証してから採用する}%
  {正解は (c)。本講義の AI Policy は「使う・検証する・区別する」の3つです。AI の利用は隠すのではなく、AI活用記録で明示します。(b) は「使う」、(d) は「検証する」、(a) は「区別する」にそれぞれ対応します。利用の透明性は評価の対象でもあります。}

